\section{Khảo sát các phương pháp chẩn đoán dựa trên ảnh}
\label{sec:methods}

\subsection{Các cách phát biểu bài toán}

Các công trình đã công bố về phân tích tự động xuất huyết não (ICH) giải quyết
nhiều bài toán khác nhau, vốn thường bị đánh đồng nhưng lại khác biệt về chi phí
gán nhãn và giá trị lâm sàng:

\begin{description}[leftmargin=2.6cm,style=nextline]
    \item[Phân loại lát cắt] gán nhãn cho từng lát cắt ngang trục. Chi phí gán
          nhãn thấp, và các tập dữ liệu công khai lớn đều được gán nhãn theo cách
          này, nhưng đầu ra cần được tổng hợp lại mới có ý nghĩa lâm sàng.
    \item[Phân loại ca chụp] gán một nhãn duy nhất cho mỗi lần thăm khám. Cách
          này khớp trực tiếp với tình huống phân loại ưu tiên, nhưng phải giải
          quyết bài toán tổng hợp: một ca chụp là dương tính nếu \emph{bất kỳ}
          lát cắt nào dương tính, nên phép lấy trung bình ngây thơ sẽ làm loãng
          các tổn thương nhỏ.
    \item[Phân vùng] khoanh vùng khối xuất huyết ở mức voxel, cho phép đo thể
          tích. Thể tích khối máu tụ là đầu vào trực tiếp cho các quyết định điều
          trị, nên đây là đầu ra có giá trị lâm sàng cao nhất, song việc gán nhãn
          ở mức voxel lại tốn kém và khan hiếm.
    \item[Phát hiện / định vị] cho ra các vùng bao, nằm ở mức trung gian cả về
          chi phí lẫn giá trị sử dụng.
\end{description}

\subsection{Các phương pháp cổ điển}

Các hướng tiếp cận trước thời kỳ học sâu khai thác đặc điểm rằng vùng xuất huyết
nằm trong một khoảng HU đặc trưng, kết hợp phép lấy ngưỡng với các phép toán hình
thái và phương pháp tăng trưởng vùng, đôi khi kèm theo một bộ phân loại truyền
thống trên các đặc trưng kết cấu được thiết kế thủ công. Các phương pháp này có
chi phí tính toán không đáng kể và hoàn toàn có thể diễn giải, đồng thời vẫn giữ
vai trò như một phép kiểm tra tính hợp lý và một phương án dự phòng. Điểm yếu của
chúng là sự thiếu bền vững: ngưỡng cố định thất bại trước hiệu ứng thể tích từng
phần ở ranh giới tổn thương, trước nhiễu ảnh do cứng hóa chùm tia gần nền sọ, và
trước vôi hóa, vốn có cùng khoảng HU với máu tụ cấp nhưng lại không có ý nghĩa
lâm sàng. Độ nhạy được báo cáo thấp hơn nhiều so với mức hiện có thể đạt được, và
không có hệ thống thuần cổ điển nào đang được sử dụng lâm sàng một cách đáng kể.

\subsection{Mạng tích chập hai chiều}

Cách phát biểu chiếm ưu thế xem mỗi lát cắt ngang trục như một ảnh và áp dụng một
mạng nơ-ron tích chập (convolutional neural network, CNN) xương sống tiêu chuẩn
(ResNet, EfficientNet, DenseNet) đã được tiền huấn luyện trên ImageNet. Phép kết
hợp nhiều cửa sổ trình bày ở Mục~\ref{sec:windowing} cung cấp ba kênh đầu vào.
Cuộc thi RSNA Intracranial Hemorrhage Detection năm 2019, vốn công bố khoảng
870{,}000 lát cắt được chuyên gia gán nhãn, đã thiết lập hướng này như một đường
cơ sở mạnh và tạo ra các quy ước kiến trúc mà phần lớn công trình sau đó tuân
theo.

Hướng tiếp cận này được hỗ trợ tốt bởi công cụ, huấn luyện được trên phần cứng
khiêm tốn, và hưởng lợi từ học chuyển giao (transfer learning). Hạn chế cố hữu
của nó là bỏ qua ngữ cảnh giữa các lát cắt. Bác sĩ chẩn đoán hình ảnh xác định
một vùng tăng tỷ trọng là xuất huyết một phần nhờ theo dõi diễn biến của nó qua
các lát cắt liền kề; một bộ phân loại theo từng lát cắt không làm được điều đó.
Điều này gây thiệt hại lớn nhất với xuất huyết dưới nhện (subarachnoid
haemorrhage, SAH) dạng mỏng, khi một lát cắt đơn lẻ có thể thực sự mơ hồ.

\subsection{Mô hình chuỗi trên đặc trưng lát cắt}

Giải pháp khắc phục tiêu chuẩn vẫn giữ một CNN 2D làm bộ trích xuất đặc trưng cho
từng lát cắt, nhưng thay việc phân loại độc lập bằng một mô hình chuỗi, thường là
LSTM hai chiều hoặc bộ mã hóa transformer, chạy trên dãy vector nhúng lát cắt đã
sắp thứ tự. Cách này khôi phục ngữ cảnh giữa các lát cắt với chi phí chỉ bằng một
phần nhỏ so với tích chập 3D đầy đủ và đã thống trị bảng xếp hạng của cuộc thi
RSNA. Nó cũng giải quyết gọn bài toán tổng hợp, vì mô hình chuỗi có thể phát ra
đồng thời dự đoán ở mức lát cắt và mức ca chụp từ một biểu diễn dùng chung. Cái
giá phải trả là một quy trình hai giai đoạn với độ phức tạp huấn luyện tương ứng
cao hơn.

\subsection{Mạng tích chập ba chiều}

Thao tác trực tiếp trên khối ảnh bằng tích chập 3D mô hình hóa ngữ cảnh không
gian một cách tự nhiên. Khó khăn nằm ở khía cạnh thực tiễn: mức tiêu thụ bộ nhớ
buộc phải hoặc giảm mẫu mạnh, làm xóa mất chính những tổn thương nhỏ quan trọng
nhất, hoặc huấn luyện theo mảnh (patch), làm tái xuất hiện giới hạn về ngữ cảnh.
Không có sẵn tiền huấn luyện trên ImageNet, nên cần nhiều dữ liệu có nhãn hơn, và
khối ảnh CT có voxel bất đẳng hướng với độ dày lát cắt thay đổi giữa các phác đồ,
đòi hỏi phải lấy mẫu lại một cách nhất quán. Do vậy, các hướng tiếp cận 3D phổ
biến hơn trong phân vùng, nơi các biến thể U-Net 3D là chuẩn mực, hơn là trong
phân loại.

\subsection{Các kiến trúc phân vùng}

U-Net và các hậu duệ của nó vẫn là kiến trúc tham chiếu, với khung nnU-Net tự
cấu hình là một lựa chọn mặc định mạnh, tự động điều chỉnh tiền xử lý và cấu trúc
mạng theo tập dữ liệu. Các kiến trúc lai bổ sung transformer mô hình hóa tốt hơn
các phụ thuộc tầm xa, đổi lại đòi hỏi nhiều dữ liệu hơn. Vì nhãn ICH ở mức voxel
khan hiếm, các công trình phân vùng đã công bố thường huấn luyện trên những tập
dữ liệu nhỏ hơn hai đến ba bậc độ lớn so với tập dữ liệu sẵn có cho bài toán phân
loại.

\subsection{Tổng hợp so sánh}

\begin{table}[H]
\centering
\small
\begin{tabularx}{\textwidth}{@{}l c c X@{}}
\toprule
\textbf{Hướng tiếp cận} & \textbf{Nhu cầu dữ liệu} & \textbf{Tính toán} & \textbf{Hạn chế chính} \\
\midrule
Lấy ngưỡng / hình thái & không & tối thiểu & Thiếu bền vững trước nhiễu ảnh và vôi hóa \\
CNN 2D theo lát cắt & trung bình & thấp & Không có ngữ cảnh giữa các lát cắt \\
CNN 2D + mô hình chuỗi & trung bình & trung bình & Độ phức tạp huấn luyện hai giai đoạn \\
CNN 3D & cao & cao & Bộ nhớ; không có tiền huấn luyện; nhạy với lấy mẫu lại \\
Họ U-Net (phân vùng) & cao (mức voxel) & cao & Chi phí gán nhãn \\
\bottomrule
\end{tabularx}
\caption{So sánh các hướng tiếp cận phân tích tự động xuất huyết não.}
\label{tab:methods}
\end{table}

\subsection{Đánh giá và bài toán tỷ lệ nền}

Các con số độ chính xác được báo cáo cần được đọc một cách thận trọng, bởi tỷ lệ
hiện mắc ICH trong quần thể chụp CT sọ não không chọn lọc là thấp. Một bộ phân
loại luôn dự đoán ``không xuất huyết'' vẫn đạt độ chính xác cao nhưng vô dụng về
mặt lâm sàng, nên độ chính xác không được dùng làm chỉ số chủ đạo. Độ nhạy, độ
đặc hiệu và diện tích dưới đường cong ROC là các thước đo phù hợp, kèm theo yêu
cầu báo cáo theo từng phân nhóm xuất huyết, vì các con số tổng hợp che giấu hiệu
năng kém đối với SAH đã bàn ở Mục~\ref{sec:clinical}.

Giá trị tiên đoán dương xứng đáng được đặc biệt chú ý, vì khác với độ nhạy và độ
đặc hiệu, nó phụ thuộc vào tỷ lệ hiện mắc. Với tỷ lệ hiện mắc điển hình ở khoa
cấp cứu, ngay cả một bộ phân loại có độ đặc hiệu rất cao vẫn sinh ra một tỷ lệ
đáng kể dương tính giả trong số các cảnh báo của nó. Nếu những cảnh báo đó đòi
hỏi hành động, tình trạng mệt mỏi vì cảnh báo phát sinh sẽ bào mòn giá trị của hệ
thống, và cân nhắc này định hình trực tiếp cách trình bày kết quả trên giao diện.

Hai cạm bẫy khác cũng rất phổ biến. \emph{Rò rỉ dữ liệu} xảy ra khi các lát cắt
của cùng một bệnh nhân xuất hiện ở cả tập huấn luyện lẫn tập kiểm thử; việc chia
tập phải được thực hiện ở mức bệnh nhân. \emph{Dịch chuyển phân phối} xảy ra khi
một mô hình huấn luyện tại một cơ sở gặp phải máy chụp và phác đồ của cơ sở khác,
và đây là nguyên nhân chính khiến hiệu năng hồi cứu cao hơn hiệu năng thực tế khi
triển khai. Cả hai đều dẫn tới yêu cầu, được phát triển ở
Mục~\ref{sec:architecture}, rằng nền tảng phải tự giám sát phân phối dữ liệu đầu
vào của chính nó sau khi triển khai.
